%% verificacao_numerica_otimiza_tx.tex
%% Verificação Numérica da Estratégia Sequencial Tx->Rx:
%% Simulação com Geometria Real, Reotimização do Receptor e
%% o Papel do Ângulo de Brewster na Equação de Radar Completa
%%
%% Idioma: Português
%% Compilar: pdflatex -interaction=nonstopmode verificacao_numerica_otimiza_tx.tex  (duas passagens)

\documentclass[12pt,a4paper,oneside]{book}

%% ─── Pacotes ─────────────────────────────────────────────────────────────────
\usepackage[utf8]{inputenc}
\usepackage[T1]{fontenc}
\usepackage[portuguese]{babel}
\usepackage{lmodern}
\usepackage{microtype}
\usepackage{geometry}
\geometry{left=3cm, right=2.5cm, top=3cm, bottom=2.5cm}

\usepackage{amsmath,amssymb,amsfonts,amsthm}
\usepackage{mathtools}
\usepackage{bm}
\usepackage{siunitx}
\sisetup{output-decimal-marker={,}}

\usepackage{graphicx}
\graphicspath{{figures/}}

\usepackage{float}
\usepackage{caption}
\usepackage{subcaption}
\usepackage[table,xcdraw]{xcolor}
\usepackage{booktabs}
\usepackage{array}
\usepackage{multirow}
\usepackage{longtable}
\usepackage{tabularx}
\usepackage{setspace}
\onehalfspacing

\usepackage{algorithm}
\usepackage{algpseudocode}

\usepackage{tcolorbox}
\tcbuselibrary{theorems,skins,breakable}

\usepackage[unicode,colorlinks=true,
  linkcolor=blue!60!black,
  citecolor=green!50!black,
  urlcolor=cyan!60!black,
  pdftitle={Verificacao Numerica da Otimizacao Sequencial Tx-Rx e o Papel do Angulo de Brewster},
  pdfauthor={David Atencia}]{hyperref}

\usepackage{fancyhdr}
\setlength{\headheight}{15pt}
\pagestyle{fancy}
\fancyhf{}
\fancyhead[L]{\small\nouppercase{\leftmark}}
\fancyhead[R]{\small\thepage}
\renewcommand{\headrulewidth}{0.4pt}

%% ─── Caixas de destaque (mesmo estilo de otimiza_tx.tex) ──────────────────────
\newcommand{\result}[2]{%
  \begin{tcolorbox}[enhanced, breakable,
    colback=yellow!8, colframe=orange!70!black,
    fonttitle=\bfseries,
    title={\protect#1},
    attach boxed title to top left={yshift=-2mm, xshift=5mm},
    boxed title style={colback=orange!70!black}]
    #2
  \end{tcolorbox}}

\newcommand{\nota}[2]{%
  \begin{tcolorbox}[enhanced, breakable,
    colback=green!5, colframe=green!50!black,
    fonttitle=\bfseries\small,
    title={\protect#1}]
    #2
  \end{tcolorbox}}

\newcommand{\atencion}[2]{%
  \begin{tcolorbox}[enhanced, breakable,
    colback=red!5, colframe=red!60!black,
    fonttitle=\bfseries\small,
    title={\protect#1}]
    #2
  \end{tcolorbox}}

%% ─── Teoremas ────────────────────────────────────────────────────────────────
\newtheorem{teorema}{Teorema}[chapter]
\newtheorem{proposicao}{Proposição}[chapter]
\newtheorem{lema}{Lema}[chapter]
\newtheorem{corolario}{Corolário}[chapter]
\theoremstyle{definition}
\newtheorem{definicaom}{Definição}[chapter]
\theoremstyle{remark}
\newtheorem*{observacao}{Observação}

%% ─── Macros ──────────────────────────────────────────────────────────────────
\newcommand{\RT}{R_T}
\newcommand{\RR}{R_R}
\newcommand{\thetaB}{\theta_B}
\newcommand{\thetaBar}{\theta_{B,\mathrm{ar}}}
\newcommand{\thetaBsolo}{\theta_{B,\mathrm{solo}}}
\newcommand{\nuno}{n_1}
\newcommand{\ndos}{n_2}
\newcommand{\Prcv}{P_r}
\newcommand{\Pt}{P_t}
\newcommand{\dxy}{\delta_{xy}}
\newcommand{\dz}{\delta_z}
\newcommand{\psiT}{\psi_{0,\mathrm{Tx}}}
\newcommand{\psiR}{\psi_{0,\mathrm{Rx}}}

% ─────────────────────────────────────────────────────────────────────────────
\begin{document}

\frontmatter

%%─── Capa ────────────────────────────────────────────────────────────────────
\begin{titlepage}
\centering
\vspace*{2cm}
{\LARGE\bfseries Verificação Numérica da Estratégia\\
Sequencial Tx$\,\to\,$Rx:\\[0.3cm]
Simulação com Geometria Real, Reotimização\\
do Receptor e o Papel do Ângulo de Brewster\\
na Equação de Radar Completa}\\[1.2cm]
{\large\itshape Resultados experimentais, diagnóstico das discrepâncias\\
frente à teoria de forma fechada, e síntese comparativa\\
de três estratégias de projeto do enlace biestático}\\[2cm]
{\normalsize Documento técnico complementar a\\
\texttt{otimiza\_tx.tex} (Otimização do Transmissor no SAR\\
Biestático Helicoidal) e à dissertação\\
\textit{Otimização da Trajetória do Receptor em SAR Biestático Helicoidal}}\\[2cm]
{\normalsize David Atencia}\\[0.4cm]
{\normalsize 2026}
\vfill
\end{titlepage}

%%─── Resumo ──────────────────────────────────────────────────────────────────
\chapter*{Resumo}
\addcontentsline{toc}{chapter}{Resumo}

O documento \texttt{otimiza\_tx.tex} propôs uma estratégia sequencial de duas
etapas para o projeto do enlace biestático -- primeiro o transmissor (Tx),
depois o receptor (Rx), ambos guiados pelo ângulo de Brewster -- e validou-a
apenas de forma analítica, sobre uma aproximação de primeira ordem
($\theta_{i,1}\approx\psi_0$, alvo único no centroide). Este documento fecha
esse ciclo: implementa a Etapa~1 e a Etapa~2 com a geometria \emph{real} do
sistema (refração de Fermat ponto a ponto, grade completa de alvos,
reotimização de 5 arranques do Rx por posição de Tx) e confronta os
resultados contra a teoria de forma fechada e contra a metodologia atual de
uma única etapa (apenas Rx otimizado).

Três achados centrais emergem. \textbf{Primeiro}, a simulação com geometria
real reproduz a \emph{forma} da curva teórica de $T_1(\psi_0)$ e
$\dz(\psi_0)$ com erro pequeno (RMS $\le 0{,}03$ em $T_1$), mas o ângulo
ótimo $\psi_0^*$ se desloca $4$--$9^\circ$ em relação à previsão teórica,
por efeito da aproximação de alvo único. \textbf{Segundo}, e mais
relevante: sob o peso de projeto atualmente em uso ($\gamma=\alpha_{\mathrm{res}}=0{,}3$),
o plano de voo de duas etapas resultante é \emph{pior} que a metodologia
atual de uma etapa em todas as quatro métricas (potência, $\dxy$, $\dz$ e
custo combinado) -- porque a demonstração de ``ganho aditivo'' do documento
teórico assume implicitamente que a posição do Rx não muda com $\psi_0$ do
Tx, o que deixa de valer quando o Rx é reotimizado de fato. Sob peso de
potência pura ($\gamma=\alpha_{\mathrm{res}}=0$), a mesma estratégia
\emph{sim} produz um ganho real ($+0{,}81\,\mathrm{dB}$, acima até do
previsto só por $T_1$). \textbf{Terceiro}, uma investigação à parte
confirma que os dois ângulos de Brewster do enlace (ar$\to$solo,
$63{,}43^\circ$; solo$\to$ar, $26{,}57^\circ$, complementares) estão
corretamente implementados em toda a base de código por três caminhos
matematicamente distintos, e que a otimização de potência pura
\emph{não} persegue Brewster -- por um motivo físico preciso e
quantificado aqui (a perda por espalhamento $1/R^2$ domina o termo de
transmitância $T(\theta)$ por mais de uma ordem de grandeza de faixa
dinâmica). Isolando artificialmente a transmitância da equação de radar,
confirma-se que o mecanismo de busca de Brewster em si está correto
(desvio médio de $0{,}04^\circ$ em 275 posições).

%%─── Sumário ─────────────────────────────────────────────────────────────────
\tableofcontents

\mainmatter

%%═══════════════════════════════════════════════════════════════════════════════
\chapter{Introdução e Objetivos desta Verificação}
%%═══════════════════════════════════════════════════════════════════════════════

\section{O que já estava estabelecido}

O documento \texttt{otimiza\_tx.tex} demonstrou, de forma analítica, que a
potência biestática recebida se fatora exatamente em um termo dependente
apenas do Tx e um termo dependente apenas do Rx (Proposição de
separabilidade multiplicativa), e propôs, a partir disso, uma estratégia
sequencial de duas etapas:

\begin{itemize}
  \item \textbf{Etapa~1 (Tx):} escolher o ângulo de \textit{look} $\psi_0$ da
  hélice do Tx minimizando $J_{\mathrm{Tx}}(\psi_0;\gamma)=-(1-\gamma)\log_{10}T_1(\psi_0)+\gamma\log_{10}\dz(\psi_0)$,
  por busca de seção áurea.
  \item \textbf{Etapa~2 (Rx):} com o Tx fixo, otimizar a posição $(x,y)$ do
  Rx a cada pulso minimizando $J_{\mathrm{Rx}}(x,y;\alpha_{\mathrm{res}})=-(1-\alpha_{\mathrm{res}})\log_{10}\Prcv+\alpha_{\mathrm{res}}\log_{10}(\dxy\dz)$
  -- a metodologia já validada e em uso.
\end{itemize}

A validação numérica apresentada ali, porém, usava exclusivamente a
aproximação de primeira ordem $\theta_{i,1}\approx\psi_0$ (alvo único no
centroide, sem refração de Fermat) para calcular $T_1$, e \textbf{não}
executava a Etapa~2 de fato -- ou seja, nunca se verificou o que acontece
quando o Rx é realmente reotimizado para o Tx que a Etapa~1 produz. O
próprio documento identificava isso como limitação e como trabalho futuro
(Seções~7.2--7.3 de \texttt{otimiza\_tx.tex}).

\section{Perguntas respondidas neste documento}

\begin{enumerate}
  \item A simulação com geometria real (refração de Fermat, grade completa
  de alvos) reproduz o $\psi_0^*$ previsto pela forma fechada? Que erro
  numérico e que fonte física explicam a diferença, se houver?
  \item Ao executar a Etapa~2 de fato sobre o Tx que a Etapa~1 produz, o
  plano de voo resultante é realmente melhor que a metodologia atual (Rx
  otimizado, Tx fixo)? Esse resultado depende do peso
  $\gamma=\alpha_{\mathrm{res}}$ escolhido?
  \item Os dois ângulos de Brewster do enlace -- ar$\to$solo e solo$\to$ar,
  que são \emph{diferentes} e complementares -- estão calculados
  corretamente em cada trecho do código de produção?
  \item Por que a otimização de potência pura não posiciona o Rx exatamente
  em Brewster, mesmo sendo esse o critério que maximiza a transmitância?
  \item Isolando a transmitância do resto da equação de radar, o mecanismo
  de busca de Brewster funciona como esperado?
\end{enumerate}

\nota{Sobre a natureza deste documento}{
Este é um relatório de \textbf{verificação numérica}, não uma proposta de
método novo. Os resultados aqui não substituem \texttt{otimiza\_tx.tex}; eles
testam suas previsões contra uma implementação de geometria real e
reportam onde a teoria se sustenta, onde se desvia, e por quê -- inclusive
quando o resultado é negativo (Capítulo~\ref{cap:pipeline_completo}).
}

%%═══════════════════════════════════════════════════════════════════════════════
\chapter{Metodologia de Simulação com Geometria Real}
\label{cap:metodologia}
%%═══════════════════════════════════════════════════════════════════════════════

\section{Diferença frente à validação analítica original}

A validação original de \texttt{otimiza\_tx.tex} calcula $T_1(\psi_0)$
avaliando o coeficiente de Fresnel TM diretamente no ângulo
$\theta_i=\psi_0$, para um único ponto (o centroide da grade de alvos), sem
resolver a refração real. Este documento substitui esse cálculo, em toda a
verificação, pelo pipeline de produção:

\begin{enumerate}
  \item Para cada posição candidata do Tx (ou do Rx), resolve-se o ponto de
  refração exato na interface $z=0$ por Newton-Raphson sobre a condição de
  Snell (3--5 iterações), não por aproximação de \textit{look angle}.
  \item O ângulo de incidência é medido entre o segmento real
  fonte$\to$refração e a normal -- no lado correto do meio (ar para o
  trecho Tx$\to$alvo; solo para o trecho alvo$\to$Rx), exatamente como no
  código de produção.
  \item $T_1$ (ou $T_2$) é calculado por posição e \textbf{promediado sobre
  toda a grade de alvos} ($4\times4\times4=64$ alvos, extensão horizontal
  $\pm40\,\mathrm{m}$, profundidade $-2$ a $-10\,\mathrm{m}$) e sobre várias
  posições ao longo da hélice candidata (8 amostras azimutais), não em um
  único ponto no centroide.
  \item $\dz$ usa a mesma fórmula $W_z$ já validada
  (Eq.~\eqref{eq:Wz_verif}), avaliada no ângulo de \textit{look} real de
  cada ponto simulado.
\end{enumerate}

\begin{equation}
  \dz = \frac{c}{2W_z}, \qquad
  W_z = \ndos B\cos\theta_{t,0} + \frac{c\,B_\perp\sin\psi_0\cos\psi_0}{\lambda_0 R_0\,\ndos\cos\theta_{t,0}}
  \label{eq:Wz_verif}
\end{equation}

\section{Parâmetros do sistema usados em todas as corridas}

\begin{table}[H]
\centering
\caption{Parâmetros físicos e geométricos, extraídos dos arquivos de
configuração do sistema em uso.}
\begin{tabular}{ll}
\toprule
Grandeza & Valor \\
\midrule
Frequência portadora $f_0$ & $400\,\mathrm{MHz}$ \\
Largura de banda $B$ & $50\,\mathrm{MHz}$ \\
Comprimento de onda $\lambda_0$ & $0{,}749\,\mathrm{m}$ \\
Índices de refração $\nuno,\ndos$ & $1{,}0\,;\,2{,}0$ \\
$\thetaBar=\arctan(\ndos/\nuno)$ & $63{,}43^\circ$ \\
$\thetaBsolo=\arctan(\nuno/\ndos)$ & $26{,}57^\circ$ \\
Comprimento de arco da hélice $B_{helix}$ & $47{,}17\,\mathrm{m}$ \\
$\psi_0$ do sistema atualmente em uso & $57{,}98^\circ\approx58^\circ$ \\
Alcance oblíquo nominal $R_0$ & $188{,}68\,\mathrm{m}$ \\
Grade de alvos & $4\times4\times4=64$, $\pm40\,\mathrm{m}$ horiz., $-2$ a $-10\,\mathrm{m}$ \\
\bottomrule
\end{tabular}
\end{table}

\section{Escala do experimento}

A varredura fina de $\psi_0$ (Etapa~1) cobre $5^\circ$ a $85^\circ$ em
passos de $1^\circ$ (81 pontos) para dois cenários de restrição de missão
($R_0$ fixo e $z_0$ fixo), cada ponto promediado sobre 8 posições da hélice
candidata e 64 alvos -- do total de \emph{simulações independentes} de
refração de Fermat por corrida completa. A Etapa~2, quando executada com a
equação de radar completa, repete a busca não-convexa de 5 arranques por
\texttt{fmincon} já validada, para cada uma de até 275 posições decimadas
do Tx.

%%═══════════════════════════════════════════════════════════════════════════════
\chapter{Etapa 1 (Tx): Simulação com Geometria Real vs. Forma Fechada}
\label{cap:etapa1_verif}
%%═══════════════════════════════════════════════════════════════════════════════

\section{Resultado}

A Tabela~\ref{tab:etapa1_verif} resume o contraste, para $\gamma=0{,}3$ (o
mesmo peso já usado para $\alpha_{\mathrm{res}}$ do lado do Rx), nos dois
cenários de restrição de missão.

\begin{table}[H]
\centering
\caption{Ótimo $\psi_0^*$ da Etapa~1: simulação com geometria real vs.
forma fechada, $\gamma=0{,}3$.}
\label{tab:etapa1_verif}
\begin{tabular}{lrrrr}
\toprule
Cenário & $\psi_0^*$ simulado & $\psi_0^*$ teórico & Diferença & RMS$(T_1)$ \\
\midrule
$R_0$ fixo & $53{,}39^\circ$ & $57{,}57^\circ$ & $-4{,}18^\circ$ & $0{,}0247$ \\
$z_0$ fixo & $38{,}96^\circ$ & $47{,}56^\circ$ & $-8{,}60^\circ$ & $0{,}0215$ \\
\bottomrule
\end{tabular}
\end{table}

\begin{figure}[H]
\centering
\includegraphics[width=0.95\textwidth]{ver_fig01_etapa1_sim_vs_teoria.jpg}
\caption{$T_1(\psi_0)$ e $\dz(\psi_0)$: simulação com geometria real
(tracejado) vs. forma fechada (sólido), para os dois cenários de
restrição. As curvas praticamente coincidem em forma; o ótimo se desloca.}
\label{fig:etapa1_sim_vs_teoria}
\end{figure}

\section{Diagnóstico}

O erro ponto a ponto em $T_1$ é pequeno (RMS $\le 0{,}025$, ou seja, as
curvas $T_1(\psi_0)$ \emph{quase coincidem} -- Figura~\ref{fig:etapa1_sim_vs_teoria}).
O deslocamento do ótimo, porém, é substancial ($4$--$9^\circ$) porque
$T_1(\psi_0)$ e $\dz(\psi_0)$ são funções largas e relativamente planas
perto de seus respectivos picos (a própria Seção~4.3.1 de
\texttt{otimiza\_tx.tex} já observa que $T_1\ge0{,}90$ em uma faixa de
quase $60^\circ$). Um viés sistemático pequeno entre a curva de um único
ponto (teoria) e a média sobre os 64 alvos reais mais 8 posições da hélice
(simulação) -- originado da extensão finita da grade e da profundidade não
nula dos alvos, que a aproximação $\theta_{i,1}\approx\psi_0$ ignora --
basta para mover o argmín de forma perceptível, mesmo com pouco impacto no
valor da função em si.

\result{Conclusão do Capítulo}{
A forma fechada de \texttt{otimiza\_tx.tex} prevê corretamente a
\emph{forma} do trade-off potência--resolução do Tx, mas \textbf{não} o
ponto exato do ótimo quando o alvo não é pontual no centroide -- confirma
numericamente a limitação que o próprio documento teórico já apontava
(Seção~7.2, ``a relação $\theta_{i,1}\approx\psi_0$ é uma aproximação de
primeira ordem''). Para dimensionamento fino, deve-se usar a simulação de
geometria real, não a forma fechada isolada.
}

%%═══════════════════════════════════════════════════════════════════════════════
\chapter{Pipeline Completo (Etapas 1+2): O Ganho Aditivo Não Se Sustenta}
\label{cap:pipeline_completo}
%%═══════════════════════════════════════════════════════════════════════════════

\section{Desenho do experimento}

Usando o $\psi_0^*=53{,}39^\circ$ (cenário $R_0$ fixo, $\gamma=0{,}3$) do
capítulo anterior, construiu-se a hélice final do Tx e executou-se a
Etapa~2 completa (5 arranques de \texttt{fmincon} por posição, guiados por
Brewster, $\alpha_{\mathrm{res}}=0{,}3$) sobre 254 posições decimadas.
Comparou-se contra uma corrida fresca da metodologia atual de uma etapa
(Tx fixo no $\psi_0\approx58^\circ$ do sistema em uso, mesmo
$\alpha_{\mathrm{res}}=0{,}3$, 268 posições) -- as duas corridas com
\emph{exatamente} os mesmos parâmetros de sistema, grade de alvos e peso de
custo, diferindo apenas em qual $\psi_0$ do Tx foi usado.

\section{Resultado}

\begin{table}[H]
\centering
\caption{Duas etapas (Tx otimizado) vs. uma etapa (Tx fixo), $\gamma=\alpha_{\mathrm{res}}=0{,}3$.}
\label{tab:pipeline_03}
\begin{tabular}{lrrr}
\toprule
Métrica & Uma etapa (base) & Duas etapas & Diferença \\
\midrule
$\psi_0$ do Tx & $57{,}98^\circ$ & $53{,}39^\circ$ & $-4{,}59^\circ$ \\
Potência média & $-33{,}97\,\mathrm{dBm}$ & $-35{,}38\,\mathrm{dBm}$ & $\mathbf{-1{,}41\,dB}$ \\
$\dxy$ médio & $26{,}08\,\mathrm{cm}$ & $27{,}82\,\mathrm{cm}$ & $\mathbf{+6{,}7\%}$ (pior) \\
$\dz$ médio & $114{,}77\,\mathrm{cm}$ & $117{,}55\,\mathrm{cm}$ & $\mathbf{+2{,}4\%}$ (pior) \\
Custo $J$ médio & $4{,}3194$ & $4{,}4299$ & pior \\
\bottomrule
\end{tabular}
\end{table}

\atencion{Resultado contraintuitivo}{
As \textbf{quatro} métricas pioram ao adicionar a Etapa~1. Isso contradiz
diretamente a Proposição de ``ganho aditivo'' de \texttt{otimiza\_tx.tex}
(Seção~5.3), que afirma que o ganho de potência ao mover $\psi_0$ é
$10\log_{10}(T_1(\psi_0^{(1)})/T_1(\psi_0^{(0)}))$,
\textbf{independente} da posição do Rx escolhida pela Etapa~2.
}

\section{Diagnóstico: por que a aditividade falha}

A demonstração da Proposição de aditividade fatora $\Prcv=K\cdot
f_{\mathrm{Tx}}(\mathrm{Tx})\cdot f_{\mathrm{Rx}}(\mathrm{Rx})$ e observa
que, para a \emph{mesma} posição de Rx, a razão de potências depende só de
$f_{\mathrm{Tx}}$. Isso é algebricamente correto -- mas na prática o Rx
\textbf{não} fica na mesma posição: a Etapa~2 o reotimiza do zero para cada
$\psi_0$ do Tx, e essa nova posição também entra no termo de resolução
$\dz$, que depende do ângulo de \textit{look} do \textbf{Rx} (não do Tx) na
fórmula $W_z$ já validada.

Mediu-se o ângulo de \textit{look} real do Rx (da posição ótima ao
centroide) em cada corrida:

\begin{table}[H]
\centering
\caption{Ângulo de \textit{look} real do Rx e linha de base tomográfica
perpendicular efetiva.}
\begin{tabular}{lrr}
\toprule
& Uma etapa (base) & Duas etapas \\
\midrule
$\psi_{0,\mathrm{Rx}}$ real, média & $34{,}68^\circ$ & $32{,}78^\circ$ \\
$B_\perp$ médio & $43{,}24\,\mathrm{m}$ & $44{,}11\,\mathrm{m}$ ($B_{helix}=47{,}17\,\mathrm{m}$) \\
\bottomrule
\end{tabular}
\end{table}

$B_\perp$ na verdade \emph{melhora} ligeiramente entre as duas corridas --
descartando um descasamento $\beta$--$\psi_{0,\mathrm{Rx}}$ como causa. A
causa real está no termo tomográfico de $W_z$, proporcional a
$\sin\psi_0\cos\psi_0=\tfrac12\sin(2\psi_0)$, que é máximo em
$\psi_0=45^\circ$: \textbf{ambos} os ângulos de Rx medidos já estão abaixo
de $45^\circ$, e o da corrida de duas etapas está \emph{ainda mais}
distante ($32{,}78^\circ$ vs.\ $34{,}68^\circ$) -- o que piora, não
melhora, esse termo.

\result{Conclusão do Capítulo}{
A Proposição de ganho aditivo é válida \emph{condicionalmente}: só se
sustenta se a posição do Rx não mudar entre as duas configurações de Tx
comparadas. Como a Etapa~2 sempre reotimiza o Rx, essa condição não se
verifica em geral -- e, para os parâmetros aqui testados
($\gamma=\alpha_{\mathrm{res}}=0{,}3$), o acoplamento residual (via
$\psi_{0,\mathrm{Rx}}$ e o termo tomográfico de $W_z$) é grande o
suficiente para \textbf{reverter} o ganho esperado. O
Capítulo~\ref{cap:limitacoes_recomendacoes} retoma isso como recomendação
de trabalho futuro.
}

%%═══════════════════════════════════════════════════════════════════════════════
\chapter{Dependência do Resultado no Peso Potência--Resolução}
\label{cap:peso_gamma}
%%═══════════════════════════════════════════════════════════════════════════════

\section{Repetição do experimento com peso de potência pura}

Repetiu-se o experimento do Capítulo~\ref{cap:pipeline_completo} com
$\gamma=\alpha_{\mathrm{res}}=0$ (potência pura, sem peso de resolução) em
ambas as etapas e em ambas as corridas (uma etapa e duas etapas), mantendo
todo o resto idêntico.

\begin{table}[H]
\centering
\caption{Duas etapas vs. uma etapa, $\gamma=\alpha_{\mathrm{res}}=0$ (potência pura).}
\label{tab:pipeline_00}
\begin{tabular}{lrrr}
\toprule
Métrica & Uma etapa (base) & Duas etapas & Diferença \\
\midrule
$\psi_0$ do Tx & $57{,}98^\circ$ & $60{,}57^\circ$ & $+2{,}59^\circ$ \\
Potência média & $-33{,}31\,\mathrm{dBm}$ & $-32{,}49\,\mathrm{dBm}$ & $\mathbf{+0{,}81\,dB}$ \\
$\dz$ médio & $134{,}59\,\mathrm{cm}$ & $132{,}37\,\mathrm{cm}$ & $-1{,}7\%$ (melhor, efeito colateral) \\
\bottomrule
\end{tabular}
\end{table}

\begin{table}[H]
\centering
\caption{Ganho de potência: previsto só por $T_1(\psi_0)$ vs.\ medido (com a
Etapa~2 completa).}
\begin{tabular}{lrr}
\toprule
Cenário de peso & Ganho previsto (só $T_1$) & Ganho medido \\
\midrule
$\gamma=\alpha_{\mathrm{res}}=0{,}3$ & $+0{,}019\,\mathrm{dB}$ & $-1{,}415\,\mathrm{dB}$ \\
$\gamma=\alpha_{\mathrm{res}}=0$    & $+0{,}019\,\mathrm{dB}$ & $+0{,}811\,\mathrm{dB}$ \\
\bottomrule
\end{tabular}
\end{table}

\section{Interpretação}

Com $\gamma=0$, a Etapa~1 move $\psi_0$ de $58^\circ$ para $60{,}57^\circ$
-- \textbf{aproximando-se} de $\thetaBar=63{,}43^\circ$ -- ao passo que com
$\gamma=0{,}3$ ela o move para $53{,}39^\circ$, \textbf{afastando-se} de
Brewster rumo ao ótimo tomográfico de $45^\circ$. No primeiro caso, o
reposicionamento do Rx na Etapa~2 \emph{amplifica} o ganho já previsto por
$T_1$ isolado ($+0{,}81\,\mathrm{dB}$ medido contra $+0{,}02\,\mathrm{dB}$
previsto); no segundo, o mesmo mecanismo o \emph{reverte} por completo.

\result{Conclusão do Capítulo}{
O benefício da Etapa~1 \textbf{não é universal} -- depende de \emph{para
que lado} do ótimo tomográfico de $45^\circ$ o peso de projeto empurra
$\psi_0$. Quando a Etapa~1 se move em direção a Brewster
($\gamma$ pequeno), o acoplamento com a Etapa~2 ajuda; quando se move
para longe de Brewster ($\gamma$ maior, priorizando resolução), o mesmo
acoplamento prejudica. Um valor único de $\gamma$ (como $0{,}3$, herdado
por analogia de $\alpha_{\mathrm{res}}$) não é uma escolha segura sem
verificar seu efeito na Etapa~2 -- o que só é possível rodando o pipeline
completo, não a Etapa~1 isolada.
}

%%═══════════════════════════════════════════════════════════════════════════════
\chapter{Verificação dos Ângulos de Brewster em Cada Trecho}
\label{cap:verificacao_angulos}
%%═══════════════════════════════════════════════════════════════════════════════

\section{Os dois ângulos de Brewster do enlace}

Para uma interface TM sem perdas, o ângulo de Brewster depende de qual lado
é o meio incidente:

\begin{equation}
  \thetaBar = \arctan\!\left(\frac{\ndos}{\nuno}\right) = 63{,}43^\circ
  \quad\text{(trecho ar$\to$solo, medido no ar)}
\end{equation}
\begin{equation}
  \thetaBsolo = \arctan\!\left(\frac{\nuno}{\ndos}\right) = 26{,}57^\circ
  \quad\text{(trecho solo$\to$ar, medido no solo)}
\end{equation}

Os dois são complementares ($\thetaBar+\thetaBsolo=90^\circ$) e descrevem o
\emph{mesmo} raio físico de Brewster visto de cada lado da interface.

\section{Três implementações, uma física}

Verificou-se que a base de código calcula esses dois ângulos corretamente
no trecho Tx$\to$alvo (ar, $\thetaBar$) e no trecho alvo$\to$Rx (solo,
$\thetaBsolo$), por \textbf{três} caminhos matematicamente distintos e
equivalentes:

\begin{enumerate}
  \item \textbf{Troca explícita de meios:} o coeficiente de Fresnel é
  chamado com a ordem dos índices invertida entre os dois trechos, e o
  ângulo de incidência é medido no lado físico correto (ar ou solo) em
  cada um -- o Brewster interno resultante é $\thetaBar$ para um trecho e
  $\thetaBsolo$ para o outro, automaticamente.
  \item \textbf{Identidade algébrica de ângulos complementares:} a
  heurística de posicionamento inicial do Rx usa uma única fórmula com
  $\tan\thetaBar$ para o trecho do ar e $1/\tan\thetaBar$ para o trecho do
  solo -- que é exatamente $\tan\thetaBsolo$, pois
  $\tan\thetaBar\cdot\tan\thetaBsolo=1$ para ângulos complementares.
  \item \textbf{Reciprocidade de Fresnel:} a implementação multicamada usa
  o ângulo do lado do ar em ambos os trechos, o que é válido porque a
  refletância (e portanto a transmitância) de Fresnel é idêntica nas duas
  direções quando avaliada em ângulos conjugados por Snell, para
  \emph{qualquer} ângulo -- não só em Brewster.
\end{enumerate}

\section{Verificação visual}

Desenvolveu-se uma ferramenta de verificação que, para um par Tx--Rx
escolhido, traça o raio real refratado (Tx$\to$ponto de refração$\to$alvo,
e alvo$\to$ponto de refração$\to$Rx) junto com um vetor de referência na
direção exata de Brewster, ancorado no mesmo ponto de refração e com o
mesmo comprimento do trecho real -- permitindo comparar visualmente a
direção real contra a ideal.

\begin{figure}[H]
\centering
\includegraphics[width=0.85\textwidth]{ver_fig02_verificacao_angulos.jpg}
\caption{Verificação visual para um par Tx--Rx: rayo real (sólido) vs.
referência de Brewster (tracejado), em ambos os trechos. A separação
angular visível no trecho Tx$\to$alvo é de $6{,}62^\circ$; no trecho
alvo$\to$Rx, de $20{,}22^\circ$ (investigado no Capítulo~\ref{cap:potencia_vs_brewster}).}
\label{fig:verificacao_angulos}
\end{figure}

\result{Conclusão do Capítulo}{
Os ângulos de Brewster de ambos os trechos estão corretamente calculados
em toda a base de código, por três estratégias de implementação distintas
mas fisicamente equivalentes. Qualquer separação angular observada entre o
raio real e a referência de Brewster (como na
Figura~\ref{fig:verificacao_angulos}) reflete uma decisão real do
otimizador, não um erro de cálculo -- questão que o próximo capítulo
resolve quantitativamente.
}

%%═══════════════════════════════════════════════════════════════════════════════
\chapter{Por Que a Potência Pura Não Persegue Brewster}
\label{cap:potencia_vs_brewster}
%%═══════════════════════════════════════════════════════════════════════════════

\section{A observação}

Ao aplicar a ferramenta de verificação visual sobre uma corrida com
$\alpha_{\mathrm{res}}=0$ (potência pura), observou-se que o Rx otimizado
fica $20{,}22^\circ$ longe de $\thetaBsolo$ -- uma diferença grande, apesar
de o objetivo ser exclusivamente potência, que Brewster deveria maximizar.

\section{Decomposição numérica}

Comparou-se a posição real do Rx ótimo contra o candidato ``Brewster puro''
(o mesmo ponto que a heurística de arranque usa como uma de suas 5
sementes, sem ruído aleatório), para a mesma posição de Tx:

\begin{table}[H]
\centering
\caption{Rx ótimo real vs.\ candidato de Brewster puro (mesma posição de Tx).}
\label{tab:brewster_vs_rango}
\begin{tabular}{lrr}
\toprule
& Rx ótimo (real) & Rx em Brewster puro \\
\midrule
Distância horizontal ao alvo & $26{,}4\,\mathrm{m}$ & $229{,}5\,\mathrm{m}$ \\
Ângulo de incidência (solo) & $6{,}35^\circ$ & $26{,}57^\circ$ (Brewster exato) \\
$T_2$ & $0{,}894$ & $1{,}000$ \\
$R_R$ (alvo$\to$Rx total) & $122{,}2\,\mathrm{m}$ & $260{,}0\,\mathrm{m}$ \\
Potência recebida & $\mathbf{-35{,}57\,dBm}$ & $-40{,}27\,\mathrm{dBm}$ \\
\bottomrule
\end{tabular}
\end{table}

O ótimo real vence por $+4{,}70\,\mathrm{dB}$. Decompondo essa diferença:

\begin{table}[H]
\centering
\begin{tabular}{lr}
\toprule
Componente & Contribuição \\
\midrule
Só por $T_2$ (transmitância) & $-0{,}48\,\mathrm{dB}$ (Brewster ganharia) \\
Só por $R_R^2$ (espalhamento geométrico) & $\mathbf{+6{,}56\,dB}$ (posição próxima ganha) \\
\bottomrule
\end{tabular}
\end{table}

\section{Explicação física}

A equação de radar biestático é $\Prcv\propto T_1T_2/(\RT^2\RR^2)$. $T_2$
só pode variar entre $\approx0{,}33$ (rasante) e $1{,}0$ (Brewster exato)
-- no máximo $\approx4{,}8\,\mathrm{dB}$ de faixa dinâmica total, e no caso
medido, como o ângulo real ($6{,}35^\circ$) já está perto de incidência
normal, $T_2=0{,}894$ não está longe do máximo (perder-se-ia só
$0{,}48\,\mathrm{dB}$ ao não estar em Brewster). Em compensação, $\RR^2$
depende do quadrado da distância: triplicar a distância para alcançar
Brewster ($26\,\mathrm{m}\to230\,\mathrm{m}$) custa
$20\log_{10}(230/26)\approx19\,\mathrm{dB}$ de perda de propagação --
muito mais do que Brewster pode compensar.

\result{Conclusão do Capítulo}{
Quando o objetivo é potência pura, aproximar-se do alvo quase sempre vence
alinhar-se com Brewster, porque o termo $1/R^2$ tem muito mais faixa
dinâmica (dezenas de dB) que $T(\theta)$ (poucos dB). O candidato de
Brewster na heurística de arranque é apenas \emph{uma} de cinco sementes
para o otimizador não-convexo -- não uma restrição -- e o otimizador a
descarta corretamente quando ela não é competitiva. Isso também explica,
retrospectivamente, a origem do ganho de $+0{,}81\,\mathrm{dB}$ do
Capítulo~\ref{cap:peso_gamma}: vem majoritariamente do efeito de alcance,
não de uma perseguição de Brewster pelo sistema.
}

%%═══════════════════════════════════════════════════════════════════════════════
\chapter{Isolando o Critério de Brewster: Otimização de Transmitância Pura}
\label{cap:brewster_puro}
%%═══════════════════════════════════════════════════════════════════════════════

\section{Desenho do experimento}

Para confirmar que o mecanismo de busca de Brewster está correto -- e que
o resultado do Capítulo~\ref{cap:potencia_vs_brewster} é uma decisão
correta do otimizador, não um defeito de implementação -- construiu-se um
pipeline de duas etapas que otimiza \textbf{apenas} $T_1$ (Etapa~1) e
\textbf{apenas} $T_2$ (Etapa~2), removendo completamente o ganho de antena
e o termo $1/R^2$ do custo (ambos continuam sendo calculados e reportados
como diagnóstico, mas não entram na otimização).

Como $T_2$ depende só do ângulo de incidência no solo -- que por sua vez
depende só da distância radial do Rx ao alvo, não do azimute -- o conjunto
de posições que maximizam $T_2$ é um \textbf{círculo completo} ao redor do
alvo, não um ponto único. Resolveu-se essa degenerescência fixando o
azimute (oposto ao Tx, mesma convenção já usada na heurística de arranque
do Rx) e buscando por seção áurea apenas a distância radial.

\section{Resultado}

\begin{table}[H]
\centering
\caption{Otimização de transmitância pura: ângulo real vs.\ Brewster, 275 posições.}
\label{tab:brewster_puro}
\begin{tabular}{lrrr}
\toprule
& Ângulo médio & $\theta_B$ correspondente & $\lvert$diferença média$\rvert$ \\
\midrule
Etapa~1 (Tx, ar) & $\psi_0^*=60{,}57^\circ$ & $\thetaBar=63{,}43^\circ$ & $2{,}87^\circ$ \\
Etapa~2 (Rx, solo, via Fermat) & $26{,}57^\circ$ & $\thetaBsolo=26{,}57^\circ$ & $\mathbf{0{,}04^\circ}$ \\
\bottomrule
\end{tabular}
\end{table}

$T_1$ médio $=0{,}984$; $T_2$ médio $=0{,}986$ -- ambos muito perto do
máximo teórico de $1{,}0$. Como diagnóstico (não otimizado): potência
média $-36{,}48\,\mathrm{dBm}$, $\dz$ médio $136{,}94\,\mathrm{cm}$ -- pior
em potência que a otimização de potência completa
(Capítulo~\ref{cap:potencia_vs_brewster}), confirmando outra vez que
Brewster puro sacrifica alcance.

\result{Conclusão do Capítulo}{
Isolada do resto da equação de radar, a Etapa~2 encontra o ângulo de
Brewster do solo com desvio médio de apenas $0{,}04^\circ$ em 275
posições -- confirma que o mecanismo de refração de Fermat, medição de
ângulo e busca por seção áurea está correto. A distância de $20{,}22^\circ$
observada no Capítulo~\ref{cap:potencia_vs_brewster} não é um erro: é o
resultado correto de uma otimização de potência completa, onde $1/R^2$
domina.
}

%%═══════════════════════════════════════════════════════════════════════════════
\chapter{Síntese Comparativa dos Três Enfoques}
\label{cap:sintese}
%%═══════════════════════════════════════════════════════════════════════════════

\section{Os três enfoques}

\begin{enumerate}
  \item \textbf{Só Rx (linha de base):} Tx fixo (hélice do sistema em
  uso), apenas o Rx otimizado pela equação de radar completa a cada pulso.
  \item \textbf{Tx+Rx, equação de radar:} Etapa~1 e Etapa~2, ambas
  otimizando a equação de radar completa (potência e/ou resolução,
  conforme $\gamma$/$\alpha_{\mathrm{res}}$).
  \item \textbf{Tx+Rx, transmitância pura:} Etapa~1 e Etapa~2, ambas
  otimizando só $T_1$/$T_2$ (Capítulo~\ref{cap:brewster_puro}).
\end{enumerate}

\begin{table}[H]
\centering
\caption{Comparação agregada dos três enfoques (corrida de potência pura,
$\gamma=\alpha_{\mathrm{res}}=0$ para os enfoques 1--2).}
\label{tab:sintese}
\begin{tabular}{lrrrr}
\toprule
Enfoque & $\psi_0$ [$^\circ$] & Potência [dBm] & $\dz$ [cm] & Ganho de potência vs.\ base \\
\midrule
Só Rx (base)          & $57{,}98$ & $-33{,}31$ & $134{,}59$ & --- \\
Tx+Rx, eq.\ radar      & $60{,}57$ & $-32{,}49$ & $132{,}37$ & $+0{,}81\,\mathrm{dB}$ \\
Tx+Rx, transm.\ pura   & $60{,}57$ & $-36{,}86$ & $136{,}94$ & $-3{,}56\,\mathrm{dB}$ \\
\bottomrule
\end{tabular}
\end{table}

\begin{figure}[H]
\centering
\includegraphics[width=0.95\textwidth]{ver_fig03_comparacao_3enfoques_3d.jpg}
\caption{Trajetórias Tx/Rx dos três enfoques sobrepostas (vista 3D e planta).
Azul: linha de base. Laranja: equação de radar completa. Verde:
transmitância pura -- o Rx traça um círculo muito mais amplo, perseguindo
Brewster sem a penalização de $1/R^2$.}
\label{fig:sintese_3d}
\end{figure}

\begin{figure}[H]
\centering
\includegraphics[width=0.85\textwidth]{ver_fig04_comparacao_3enfoques_barras.jpg}
\caption{Potência média, $\dz$ médio e $\psi_0$ do Tx, com desvio padrão,
para os três enfoques.}
\label{fig:sintese_barras}
\end{figure}

\result{Conclusão do Capítulo}{
Entre os três, a equação de radar completa (enfoque~2) é a única que supera
a linha de base nesta corrida específica -- e por uma margem modesta
($+0{,}81\,\mathrm{dB}$), majoritariamente por efeito de alcance
(Capítulo~\ref{cap:potencia_vs_brewster}), não por perseguição de Brewster.
A transmitância pura (enfoque~3) perde potência de forma considerável
frente aos outros dois, por desenhar deliberadamente um enlace mais longo
para maximizar $T$ -- é a estratégia certa \emph{se e somente se}
transmitância for de fato a métrica que importa, não potência recebida.
}

%%═══════════════════════════════════════════════════════════════════════════════
\chapter{Conclusões, Limitações e Recomendações}
\label{cap:limitacoes_recomendacoes}
%%═══════════════════════════════════════════════════════════════════════════════

\section{Síntese dos achados}

\begin{enumerate}
  \item A simulação com geometria real confirma a forma da teoria de
  \texttt{otimiza\_tx.tex} para a Etapa~1 (Capítulo~\ref{cap:etapa1_verif}),
  mas o ótimo $\psi_0^*$ se desloca $4$--$9^\circ$ por efeito da
  aproximação de alvo único -- limitação já antecipada no documento
  teórico, agora quantificada.
  \item A Proposição de ganho aditivo (Etapa~1 independente da Etapa~2)
  \textbf{não se sustenta} quando o Rx é de fato reotimizado: sob o peso
  de projeto atual ($\gamma=\alpha_{\mathrm{res}}=0{,}3$), o pipeline de
  duas etapas piora as quatro métricas frente à metodologia de uma etapa
  (Capítulo~\ref{cap:pipeline_completo}).
  \item Esse resultado \textbf{depende do sinal} do deslocamento de
  $\psi_0$ relativo a $45^\circ$: sob potência pura
  ($\gamma=\alpha_{\mathrm{res}}=0$), o mesmo pipeline \emph{ganha}
  $+0{,}81\,\mathrm{dB}$ (Capítulo~\ref{cap:peso_gamma}).
  \item Os dois ângulos de Brewster do enlace estão corretamente
  implementados em toda a base de código (Capítulo~\ref{cap:verificacao_angulos}).
  \item A otimização de potência pura não persegue Brewster porque o termo
  $1/R^2$ domina $T(\theta)$ em faixa dinâmica -- confirmado numericamente
  ($+6{,}56\,\mathrm{dB}$ de alcance vs.\ $-0{,}48\,\mathrm{dB}$ de
  transmitância, Capítulo~\ref{cap:potencia_vs_brewster}) e por
  contraprova isolando a transmitância pura
  (Capítulo~\ref{cap:brewster_puro}, desvio de $0{,}04^\circ$).
\end{enumerate}

\section{Limitações}

\begin{itemize}
  \item Os resultados dos Capítulos~\ref{cap:pipeline_completo}
  e~\ref{cap:peso_gamma} foram obtidos para o cenário $R_0$ fixo e um
  único par de valores de $\gamma=\alpha_{\mathrm{res}}$ de cada vez; não
  se varreu $\gamma$ continuamente para mapear onde exatamente o sinal do
  ganho se inverte, nem se repetiu para o cenário $z_0$ fixo.
  \item A comparação de três enfoques do Capítulo~\ref{cap:sintese} usa
  corridas específicas de potência pura para os enfoques 1--2; os valores
  numéricos mudariam sob outro peso de projeto (como já demonstra o
  Capítulo~\ref{cap:peso_gamma} para os enfoques 1--2 isolados).
  \item A distância de Brewster medida no Capítulo~\ref{cap:potencia_vs_brewster}
  foi decomposta para uma única posição de Tx representativa; não se
  verificou se a decomposição $T$ vs.\ $1/R^2$ muda de sinal em algum
  ponto da hélice.
\end{itemize}

\section{Recomendações}

\begin{itemize}
  \item \textbf{Não adotar $\gamma=\alpha_{\mathrm{res}}$ por analogia.} A
  Etapa~1 deveria ser calibrada rodando o pipeline completo (Etapas~1+2)
  para uma grade de valores de $\gamma$, escolhendo o que de fato melhora
  as métricas finais -- não copiando o peso já usado do lado do Rx.
  \item \textbf{Revisar a Proposição de ganho aditivo} de
  \texttt{otimiza\_tx.tex} para incorporar explicitamente o acoplamento via
  $\psi_{0,\mathrm{Rx}}$ real (não assumir posição de Rx fixa), ou
  reformulá-la como um limite superior válido apenas quando o Rx não muda.
  \item \textbf{Formalizar o critério de decisão} entre otimizar potência
  vs.\ transmitância: são objetivos genuinamente distintos
  (Capítulo~\ref{cap:sintese}), e a escolha deveria ser explícita no
  projeto do sistema, não implícita na escolha de $\alpha_{\mathrm{res}}$.
  \item Estender esta verificação com uma varredura fina de $\gamma$ (como
  já existe para a Etapa~1 isolada) executando a Etapa~2 completa em cada
  ponto, para mapear a curva de ganho real vs.\ $\gamma$ e identificar o
  ponto de cruzamento onde o sinal se inverte.
\end{itemize}

\end{document}
